% ===== 原文 pp.16–18（§3.1 Teichmüller 分支收尾；§3.2 正则纤维） =====

% [原文 p.16]
\emph{Teichmüller 分支} [62]，即所谓的 Hitchin 纤维化的 \emph{Hitchin 截面}，由 $\fg_{\C}$ 的
Kostant 切片 $\mathcal{K}$ 诱导，后者是由
\[
\mathcal{K}:=\{f\in\fg_{\C}\mid f=f_0+a_1e_1+a_2e_2+\ldots+a_re_r\}
\]
给出的空间，其中 $e_i$ 是 $\fg_{\C}$ 的最高权向量，$a_i$ 是复数，而 $f_0$ 是 $\fg_{\C}$ 的一个
三维子代数 $\mathfrak{s}$ 中的幂零元。回忆：此子代数由半单元素 $h_0$ 与 $\fg_{\C}$ 中的幂零元素
$e_0$、$f_0$ 生成，它们满足关系
\[
[h_0,e_0]=e_0\,;\quad [h_0,f_0]=-f_0\,;\quad [e_0,f_0]=h_0.
\]
% [Tutor Note: p.16 这里的 $k_i$ 是 $\mathfrak{g}_{\C}$ 的指数（exponents），而式 (3.2) 与 Hitchin 底
%  定义中 $p_i$ 的次数 $d_i$ 满足 $d_i=k_i+1$；正因如此，下文才有 $a_i\in H^0(\Sigma,K^{k_i+1})$ 与
%  $\mathcal{A}_{G_{\C}}=\bigoplus_i H^0(\Sigma,K^{d_i})$ 的相容（$p_i(\Phi)=a_i$）。可用
%  $G_{\C}=\Sl(2,\C)$ 验算：唯一指数 $k_1=1$，故 $d_1=2$，$a_1=\omega$ 是二次微分。]

由 $\fg_{\C}$ 关于 $\mathrm{ad}_{h_0}$ 的特征空间分解，考虑向量丛
\[
E=\mathrm{ad}P\otimes\C=\bigoplus_{j=-k_r}^{k_r}\mathfrak{g}_j\otimes K^j,
\]
其中 $k_r=\max\{k_i\}_i$ 是 $\fg_{\C}$ 的最大指数。这是 $\fg_{\C}$ 的伴随丛，关联于主
$\mathrm{Ad}G_{\C}$-丛 $P=P_1\times_i G$，其中 $P_1$ 是关联于 $K^{-1/2}\oplus K^{1/2}$ 的全纯主
$\Sl(2,\C)$-丛，而包含映射 $i:\Sl(2,\C)\hookrightarrow\mathrm{Ad}G_{\C}$ 对应于主三维子代数
$\mathfrak{s}_0$。尽管 $P_1$ 涉及平方根 $K^{1/2}$ 的选取，丛 $E$ 却与此选取无关。当
\[
\Phi=f_0+a_1e_1+a_2e_2+\ldots+a_re_r,
\]
且 $a_i\in H^0(\Sigma,K^{k_i+1})$ 时，$\mathrm{Ad}G_{\C}$-Higgs 丛 $(E,\Phi)$ 是稳定的。注意
$(\mathrm{ad}h_0)f_0=-f_0$，因此可以把 $f_0$ 看作 $(\mathfrak{g}_{-1}\otimes K^{-1})\otimes K$ 的
截面。此外，最高权向量 $e_j\in\mathfrak{g}_{k_j}$，于是 $a_je_j$ 是 $\mathfrak{g}_{k_j}\otimes
K^{k_j+1}$ 的截面，这使得 $\Phi$ 成为 $E\otimes K$ 的一个良定义的全纯截面。

由于 $p_i(\Phi)=a_i$，上述构造定义了 $\mathcal{M}_{\mathrm{Ad}G_{\C}}$ 的一个截面
\begin{equation}
s:(a_1,\ldots,a_r)\mapsto(E,\Phi) \tag{3.3}
\end{equation}
其像即为 Teichmüller 分支。该分支在 $h$ 的光滑纤维中定义了一个原点，从而 Hitchin 纤维化成为
阿贝尔簇的纤维化（更多细节请读者参阅 [61]，Hitchin 正是在该文中引入了 Teichmüller 分支）下面的
结果（[62, Theorem 7.5]）把 Teichmüller 分支与 $\mathrm{Ad}G_{\C}$ 的分裂实形式的表示空间联系
起来：
% [Tutor Note: p.16 原文 "(for more details ... introduced the Teichmüller component) The following
%  result" 中，右括号之后缺句号；中译按原样保留该处断句（不补句号），仅在此说明。]

\begin{thm}{3.4}
$\mathcal{M}_{\mathrm{Ad}G_{\C}}$ 的截面 $s$ 定义了 $\pi_1(\Sigma)$ 到 $\mathrm{Ad}G_{\C}$ 的
分裂实形式中的约化表示模空间的一个光滑连通分支。
\end{thm}

\subsection{Hitchin 纤维化的正则纤维}

我们首先考虑 Hitchin 纤维化的正则纤维，它们是 Hitchin 底中正则的——即一般的——点上的纤维。
确切地说，我们将考虑 Hitchin 底中那些定义出\emph{光滑}谱曲线 $S$ 的点
$a\in\mathcal{A}_{G_{\C}}$ 上的纤维。如前，取 $K=T^*\Sigma$，

% [原文 p.17]
$X=\mathrm{Tot}(K)$ 为其全空间，投影为 $\pi:X\to\Sigma$。我们用 $\eta$ 记 $X$ 上拉回丛
$\pi^*K$ 的重言截面（tautological section）。我们滥用记号，用同一符号表示 $\Sigma$ 上 $K^i$ 的幂次丛的截面及其到 $X$
的拉回（更多细节见 [18] 中的方程，那里在定义谱数据时作了区分）。

一般纤维 $h^{-1}(a)$ 中的 Higgs 丛 $(E,\Phi)$ 的特征多项式在 $X$ 中定义了一条光滑曲线
$\pi:S_a\to\Sigma$，即 $\Phi$ 的\emph{谱曲线}，其方程为
\begin{equation}
\det(\eta-\Phi)=\eta^n+a_1\eta^{n-1}+a_2\eta^{n-2}+\ldots+a_{n-1}\eta+a_n=0, \tag{3.5}
\end{equation}
其中 $a_i\in H^0(\Sigma,K^i)$（为简便起见，我们将略去 $S_a$ 的下标 $a$）。由 $X$ 上的伴随
公式（例如见 [81]），$\Sigma$ 的余切丛是辛流形，因而典范丛平凡，故 $K_S\otimes\pi^*K^{-m}$ 平凡，
且 $K_S=\pi^*K^n$。两边取次数即得亏格
\[
g_S=1+n^2(g-1).
\]
% [Tutor Note: p.17 原文 "so K_S ⊗ π*K^{−m} is trivial and K_S = π*K^n" 中的 m 应为 n：$S\subset X$
%  由 $\pi^*K^n$ 的截面定义，而 $K_X$ 平凡，由伴随公式 $K_S=(K_X\otimes\mathcal{O}(S))|_S$ 直接得
%  $K_S\cong\pi^*K^n$；取 m=n 才与紧随其后的 $K_S=\pi^*K^n$ 自洽。正文照原文排 "−m"。]
% [Tutor Note: 亏格公式可验算：$\deg K_S=2g_S-2$ 而 $\deg\pi^*K^n=n(2g-2)$，故
%  $2g_S-2=n(2g-2)$，即 $g_S=1+n^2(g-1)$，与原文一致（n=2 时 $g_S=4g-3$）。]

关联于 Higgs 丛 $(E,\Phi)$ 的\emph{谱数据}给出了 Hitchin 纤维化的纤维作为阿贝尔簇的几何描述，
它由以下部分给出：
\begin{itemize}[leftmargin=1.4em,itemsep=1pt,topsep=1pt]
\item 谱曲线 $\pi:S\to\Sigma$，以及
\item $S$ 上的一个向量丛，满足反映 $G_{\C}$ 性质的适当条件。
\end{itemize}

就经典 Higgs 丛而言，光滑纤维是 $\Jac(S)$。对于 $\pi^*K$ 的重言截面 $\eta$，可以由曲线
$(S,L\in\Jac(S))$ 恢复出 $(E,\Phi)$：取 $E=\pi_*L$，$\Phi=\pi_*\eta$。事实上，回忆正像层的
定义：给定开集 $\mathcal{U}\subset\Sigma$，有
$H^0(\pi^{-1}(\mathcal{U}),L)=H^0(\mathcal{U},\pi_*L)$，而乘以 $\eta$ 诱导出映射
\[
H^0(\pi^{-1}(\mathcal{U}),L)\xrightarrow{\ \eta\ }H^0(\pi^{-1}(\mathcal{U}),L\otimes\pi^*K).
\]
考虑该映射的正像，便得到
\[
\Phi:\pi_*L\to\pi_*L\otimes K,
\]
从而给出 Higgs 场 $\Phi\in H^0(\Sigma,\End E\otimes K)$，其中 $E:=\pi_*L$。此外，该 Higgs 场
满足其特征方程，按构造它由方程
$\eta^n+a_1\eta^{n-1}+a_2\eta^{n-2}+\ldots+a_{n-1}\eta+a_n=0$ 给出。又因 $S$ 不可约，由注 2.4
知该 Higgs 场不存在不变子丛，所以诱导出的 Higgs 丛 $(E,\Phi)$ 是稳定的。反之，设 $(E,\Phi)$ 是
经典 Higgs 丛。其特征多项式由式 (3.5) 给出，其系数在全空间 $X$ 中定义了谱曲线 $S$。

由 [18, Proposition 3.6]，Higgs 丛 $(E,\Phi)$ 与前述谱曲线 $S$ 上的线丛 $L$ 之间存在双射对应。
这一对应把 Hitchin 映射的纤维与适当次数的线丛构成的 Picard 簇等同起来。通过把线丛 $L$ 与一个
取定的次数为 $-\deg(L)$ 的线丛作张量积，便得到 $\Jac(S)$ 中的一点，

% [原文 p.18]
即 $S$ 上次数为零的线丛构成的阿贝尔簇，其维数为 $g_S$，如 (3.2) 所述。特别地，雅可比簇是
Picard 群 $H^1(S,\mathcal{O}^*_S)$ 中单位元所在的连通分支。于是，经典 Hitchin 纤维化
$h:\mathcal{M}\to\mathcal{A}$ 的纤维同构于谱曲线 $S$ 的雅可比簇。更多细节，读者可参阅例如
[63, Section 2]。
% [Tutor Note: p.18 首句 "which has dimension g_S as in (3.2)" 中 "(3.2)" 与内容不符：$g_S$ 的公式
%  $g_S=1+n^2(g-1)$ 是 p.17 的无编号行间式，而 (3.2) 是注 3.1 中的特征多项式。正文照原文排 (3.2)。]

\begin{exmp}{3.6}
就秩 2 的经典 Higgs 丛 $(E,\Phi)$ 而言，$\Phi$ 的特征多项式定义了一条谱曲线
$\pi:S\to\Sigma$。这是 $K$ 的全空间中 $\Sigma$ 的一个二重覆盖，其方程为 $\eta^2+a_2=0$，其中
$a_2$ 是一个二次微分，$\eta$ 是 $\pi^*K$ 的重言截面。由 [18, Remark 3.5]，当 $a_2$ 有单零点时
该曲线光滑；此时分歧点由 $a_2$ 的除子给出。设 $z$ 是分歧点附近的局部坐标，则该覆盖由
$z\mapsto z^2:=w$ 给出。在 $z=0$ 的邻域中，线丛 $M$ 的一个截面形如
$f(w)=f_0(w)+zf_1(w)$。由于 Higgs 场由乘以 $\eta$ 得到，故有
\[
\Phi\big(f_0(w)+zf_1(w)\big)=wf_1(w)+zf_0(w),
\]
从而 Higgs 场 $\Phi$ 的一个局部形式为
\[
\Phi=\begin{pmatrix} 0 & w\\ 1 & 0\end{pmatrix}.
\]
\end{exmp}
% [Tutor Note: p.18 例 3.6 的局部形式可验算：把截面按基底 $\{1,z\}$ 写成分量 $(f_0,f_1)$，则
%  "乘以 $\eta$" 的作用是 $(f_0,f_1)\mapsto(wf_1,f_0)$——这正是 $2\times2$ 矩阵
%  $\left(\begin{smallmatrix}0&w\\1&0\end{smallmatrix}\right)$ 的乘法，与本段给出的
%  $\Phi(f_0+zf_1)=wf_1+zf_0$ 完全一致，也正是例 2.5 中 $\Phi_\omega$ 的局部形状。]

通过考虑对 Higgs 场 $\Sigma$ 的额外条件（例如要求它具有如例 2.14 中的 $\U(p,q)$-结构），特征
多项式以及由它定义的曲线会获得进一步的结构，下面将描述这类例子。因此，当
$G_{\C}\subset\Gl(n,\C)$ 时，$G_{\C}$-Higgs 丛的谱数据由该对作为经典 Higgs 丛的谱数据给出，
并满足额外条件。一个例子是：当 $G$ 是复李群 $G_{\C}$ 的分裂实形式时，实 $G$-Higgs 丛的谱数据
可以如何表达。
% [Tutor Note: p.18 原文 "additional conditions on the Higgs field Σ" 中把 Σ 说成 "Higgs field" 显然
%  有误——Σ 是底黎曼面，条件应加在 Higgs 丛（或 Higgs 场 Φ 所在的丛）上；正文照原文排 "Higgs 场 Σ"。]

\begin{thm}{3.7（[89] 定理 4.12）}
当 $G$ 是复李群 $G_{\C}$ 的分裂实形式时，$G$-Higgs 丛的谱数据由 $G_{\C}$ 的 二阶谱数据给出。
换言之，Higgs 丛模空间 $\mathcal{M}_{G_{\C}}$ 中对应于 $\fg_{\C}$ 的分裂实形式的子空间与
Hitchin 纤维化
\[
h:\mathcal{M}_{G_{\C}}\to\mathcal{A}_{G_{\C}}
\]
的光滑纤维之交，由这些纤维中的 二阶元素给出。
\end{thm}
% [Tutor Note: 定理 3.7 的 "order two" 指 Prym 纤维中的 2-挠点（与 p.20 的 $P[2]$ 同一对象）：对
%  分裂实形式，谱数据满足 $\sigma^*L\cong L^*$ 且 $L^{\otimes2}\cong\mathcal{O}$，故纤维中的 2 阶元
%  集就是 $P[2]$；这也解释了为何 p.20 要用 $H^1(\Prym,\Z)$（基本群 $\wedge$ 的秩为 $2(g_S-g)$，模 2 后即 $P[2]$）来描述单值性。]

通过上述定理，可以把某些实 Higgs 丛看作 Hitchin 纤维化中的覆盖空间，从而可以通过覆盖空间的
性质来研究它们的几何——例如，通过该纤维化的单值作用。为定义 Hitchin 纤维化的单值作用，考虑
一个一般的纤维化 $p:Y\to B$，它是局部平凡的，即对任意点 $b\in B$ 都存在开邻域
$U_b\in B$ 使得 $p^{-1}(U_b)\cong U_b\times Y_b$，其中 $Y_b$
% [Tutor Note: p.18 末 "one can certain real Higgs bundles as a covering space" 原文缺动词（应为
%  "one can regard/consider certain real Higgs bundles as ..."）；中译按 "把……看作……" 补足该动词。]
